% Source: https://artofproblemsolving.com/community/c3h1243454_mc_nats_2016
% Stable identity: aops-topic-1243454-q29
Let the square $ACEG$ have side $a$. The octagon is this square plus four congruent isosceles triangles such as $ABC$ (since $AB=BC$), each with base $a$ and height $h$, so its area is $a^2+2ah$.

\begin{asy}
usepackage("amsmath");
size(4cm);
defaultpen(fontsize(10pt));
real a=7,h=1; pair[] Sq={(a/2,a/2),(-a/2,a/2),(-a/2,-a/2),(a/2,-a/2)};
pair[] Mid={(0,a/2+h),(-a/2-h,0),(0,-a/2-h),(a/2+h,0)};
draw(Sq[0]--Mid[0]--Sq[1]--Mid[1]--Sq[2]--Mid[2]--Sq[3]--Mid[3]--cycle,linewidth(1));
draw(Sq[0]--Sq[1]--Sq[2]--Sq[3]--cycle,blue); draw(Mid[0]--Mid[1]--Mid[2]--Mid[3]--cycle,red+dashed);
label("$A$",Sq[0],NE);label("$B$",Mid[0],N);label("$C$",Sq[1],NW);label("$D$",Mid[1],W);label("$E$",Sq[2],SW);label("$F$",Mid[2],S);label("$G$",Sq[3],SE);label("$H$",Mid[3],E);
label("\mbox{}",(min(currentpicture,user=true)+max(currentpicture,user=true))/2,black); // keeps the text after the figure black when LaTeX typesets the labels inline
\end{asy}

From $a^2=\dfrac{7}{9}\left(a^2+2ah\right)$: $9a^2=7a^2+14ah$, so $h=\dfrac {a}{7}$. By symmetry $BDFH$ is a square centred at the centre $O$, with $OB=\dfrac {a}{2}+h=\dfrac{9a}{14}$, so its area is $2\cdot OB^2=\dfrac{81a^2}{98}$. The octagon's area is $a^2+\dfrac{2a^2}{7}=\dfrac{9a^2}{7}$. The ratio is
\[\dfrac{81}{98}\cdot\dfrac{7}{9}=\boxed{\dfrac{9}{14}}.\]
